\KOMAoption{paper}{a4, portrait}
\areaset{267mm}{180mm}
\fancyheadoffset{0pt} % обновляем хедер и футер

\newgeometry{
  top=17mm,
  bottom=10mm,
  right=8mm,
  left=20mm,
}

\begin{landscape}

\setlength{\headsep}{7.5mm}
\setlength{\topmargin}{-27mm}
\setlength{\footskip}{21pt}

\color{uniblue}{\section[(обязательное) Подключение дискретных входов и выходов устройства к внутренним цепям алгоритмов для архитектуры II типа]{(обязательное)\\Подключение дискретных входов и выходов устройства к внутренним цепям алгоритмов для архитектуры II типа}\label{app:inouts2}}
\color{black}

%\color{uniblue}\subsection{Общие положения}
%\color{black}
%\textcolor{red}{На рисунках \ref{fig:sig1}, \ref{fig:sig2} показано типовое подключение выходов алгоритмов функций к дискретным входам и выходам в составе устройства для исполнения ИЭУ1 в составе корпоративного шкафа ШЭТ (ПАО Россети), обозначение групп оперативных токов взято из стандарта СТО Россетей}

\begin{figure}[h!]
  \centering
  \hspace*{-1cm} % Сдвигаем рисунок влево на 2 см  
  
\includegraphics[width=1.5\textwidth]{\apppath/Приложение. Подкл вх и вых2/_latex/img/1.pdf}
  \caption{Подключение дискретных входов и выходов к внутренним сигналам алгоритмов для платы в слоте М1}
  \label{fig:sig21} % Опционально: ссылка на рисунок
\end{figure}

\begin{figure}[h!]
  \centering
  %\hspace*{-1cm} % Сдвигаем рисунок влево на 2 см    
  \includegraphics[width=1.5\textwidth]{\apppath/Приложение. Подкл вх и вых2/_latex/img/2.pdf}
  \caption{Подключение дискретных входов и выходов к внутренним сигналам алгоритмов для плат в слотах М3, М4}
  \label{fig:sig22} % Опционально: ссылка на рисунок
\end{figure}

\begin{figure}[h!]
  \centering
  %\hspace*{-1.5cm} % Сдвигаем рисунок влево на 2 см    
  \includegraphics[width=1.45\textwidth]{\apppath/Приложение. Подкл вх и вых2/_latex/img/3.pdf}
  \caption{Подключение дискретных входов и выходов к внутренним сигналам алгоритмов для плат в слотах М5-М7}
  \label{fig:sig23} % Опционально: ссылка на рисунок
\end{figure}

\clearpage
\end{landscape}
\restoregeometry


